% Workstream xs: measurement of the favourite--longshot bias in the cross-section of Kalshi markets.
% Code: src/albert/research/flb_xs.py (steps scan, sample, events, candles, panel, analyze).
% Every number in this subsection's text is a design parameter of that code, not an estimate.

\paragraph{Observations.}
For market $m$ with scheduled time $T_m$ (defined in \cref{sec:data-xs}) and lead $L$, the
observation time is $t=T_m-L$. The pre-event leads are $L\in\{1,6,24,72,168,720\}$ hours. Sports
events whose ticker carries the start time also have in-play observations at $t=T_m+0.5$, $+1$
and $+2$ hours. The quote at $t$ is the YES bid $b$ and ask $a$ at the close of the last hourly
candle that ended at or before $t$. Kalshi omits many hours without activity, and the last close is taken
as the standing quote. Its age is recorded and used in a robustness check. The mid is $p=(a+b)/2$.
An observation enters the analysis when the market is open at $t$, stays open at least 15 minutes
after $t$, and its quote is \emph{usable}: $a>b$ and $a-b\le 10\cents$. A one-sided book at the
tick floor ($b=0$, $a=1\cents$) is usable and has mid $0.5\cents$.

\paragraph{Sides, prices and returns.}
Each quote yields two side observations: YES, with side price $\q=p$ and payoff $\won=y$, and NO,
with $\q=1-p$ and $\won=1-y$, where $y=1$ if the market settled YES. The executable taker price is
$\q^{a}=a$ for YES and $1-b$ for NO; a resting (maker) order buys at $\q^{b}=b$ or $1-a$. Returns
per dollar are $\ret=(\won-\q-\fee)/(\q+\fee)$ as in \cref{sec:framework}, with the taker fee
$\fee(\q)=\theta\,\q(1-\q)$, where $\theta=0.07$ times the series fee multiplier (one for almost all
series). We report three versions:
gross at the mid ($\q=p$ or $1-p$, $\fee=0$), gross at the ask ($\q=\q^a$, $\fee=0$), and net of the
exact taker fee at the ask ($\q=\q^a$, $\fee=\fee(\q^a)$). The difference between the first two is the
half-spread; the difference between the last two is the fee. Kalshi computes fees per fill to
\$0.000001 and rebates rounding over-payment \citep{kalshirounding}, so we use exact fees and show
per-order rounding to the cent only for comparison. Fees follow the schedule and
series multipliers in force in November 2026 throughout; Kalshi's schedule differed in parts of the
earlier sample, so net returns before 2025 are indicative. A side is a
\emph{longshot} (\LS) if $\q\le 10\cents$ and a \emph{favourite} (\FAV) if $\q\ge 90\cents$. Price
buckets are those of the shared definitions: 1--2, 3--5, 6--10, 11--20, 21--35, 36--50, 51--65, 66--80,
81--90, 91--95, 96--98\cents{} and 99\cents.

\paragraph{Weights.}
A market has up to six pre-event observations. \emph{Market-weighted} statistics give each
observation weight $1/n_m$, where $n_m$ is the number of observations of market $m$ in the sample
being summarised, so that every market counts once. \emph{Volume-weighted} statistics multiply this
by the number of contracts traded in market $m$ in the seven days up to $t$. That count is known at
$t$. We do not weight by lifetime volume, because it includes trading after $t$, and a longshot
that later wins attracts heavy trading on its way up.

\paragraph{Three measures of the bias.}
\begin{enumerate}
\item \emph{Returns by price bucket}: the weighted mean of $\ret$ within each bucket of the side price
      (the mid for gross returns, the taker price for net returns).
\item \emph{Calibration slope}: the logit regression of \cref{sec:framework}, estimated with weights
      on YES sides,
      \begin{equation}\label{eq:xs-logit}
        \Prob(y=1\mid p)=\Lambda\bigl(\alpha+\beta\,\operatorname{logit}p\bigr),
      \end{equation}
      with $p$ clipped to $[0.5\cents,99.5\cents]$. Calibrated prices have $\alpha=0$ and $\beta=1$.
      With $\beta>1$, true odds are more extreme than quoted odds: longshots are overpriced and
      favourites underpriced, which is the favourite--longshot bias. Fitting on NO sides gives the
      same $\beta$ and the opposite $\alpha$, so $\beta$ summarises both tails and $\alpha$ is the
      YES-only intercept (a preference for buying YES appears as $\alpha<0$). We report
      $t_{\beta=1}=(\hat\beta-1)/\mathrm{se}(\hat\beta)$.
\item \emph{Longshot--favourite spread}: the difference between the mean return per dollar of
      \LS{} sides and of \FAV{} sides. We estimate it as the coefficient on a longshot indicator in a
      weighted regression of $\ret$ on a constant and that indicator, over \LS{} and \FAV{} sides.
\end{enumerate}
Standard errors are clustered by event \citep{cameronmiller2015}, because the legs of an event and the
leads of a market share one outcome. Each table reports the number of observations, the number of
clusters and $t$-statistics.

\paragraph{Cross-sections.}
\begin{itemize}
\item \emph{Category}: the coarse groups of the shared definitions.
\item \emph{Horizon}: each lead $L$ separately, which gives one observation per market and lead.
\item \emph{Contract structure}, from the event and leg design:
  \begin{itemize}
  \item a \emph{ladder leg} belongs to a set of at least three legs of one event that differ only by a
        numeric strike (temperature buckets, run totals, point spreads, player-prop thresholds);
  \item a \emph{deadline} contract pays if an occurrence happens before a date $T$ stated in the
        rules, classified from the rules text with the hazard study's parser;
  \item a \emph{multi-outcome} leg belongs to an event with at least three non-ladder legs that
        Kalshi flags as mutually exclusive;
  \item everything else is a \emph{binary}: single-market events, two-leg head-to-head events and
        independent legs that share an event, such as the words of a mentions event.
  \end{itemize}
\item \emph{Liquidity}: within-group terciles of the 7-day volume, the quoted spread and the open
      interest at $t$.
\item \emph{Calendar time}: quarters of $T_m$.
\end{itemize}

\paragraph{Multi-outcome events and the Shin model.}
For a mutually exclusive event with $K$ legs, all quoted at $t$ and exactly one settling YES, the
overround is $\sum_k a_k-1$ at the asks and $\sum_k p_k-1$ at the mids. Following
\citet{shin1991,shin1993} and \cref{sec:framework}, quoted prices $\pi_k$ with $\Pi=\sum_k\pi_k>1$
map to probabilities
\begin{equation}\label{eq:xs-shin}
  \hat p_k(\shinz)=\frac{\sqrt{\shinz^{2}+4(1-\shinz)\pi_k^{2}/\Pi}-\shinz}{2(1-\shinz)},
\end{equation}
where $\shinz$ is the share of insider trading and solves $\sum_k\hat p_k(\shinz)=1$. We compute
$\shinz$ per event and lead from the asks. We then compare the realised frequency of each leg with
three probability forecasts: the raw mid, the mid normalised to sum to one, and $\hat p_k$
\citep{strumbelj2014}. Legs are grouped by their price rank within the event. We also measure the
\emph{field}: the summed price of all legs at or below 10\cents{} against the number of them that win.

\paragraph{Limits to arbitrage.}
\begin{itemize}
\item \emph{Fees.} We show the taker fee per dollar staked, $\fee(\q)/\q=\theta(1-\q)$, with and
      without per-order rounding.
\item \emph{Taking the other side.} For every quote with a longshot side, we compute the net return
      of buying the opposite (favourite) side as a taker at $\q^a$. We also compute it for resting a
      bid at $\q^b$ as a maker, paying the maker fee $\theta_m\q(1-\q)$, with $\theta_m=0.0175$ times the
      multiplier, only on series that charge it. The maker figure assumes a fill at the quote without adverse selection, so it is an upper
      bound; who actually fills is the subject of \cref{sec:method-flow}. Exploiting overpriced
      longshots this way is the arbitrage whose limits \citet{shleifervishny1997} describe.
\item \emph{No-arbitrage band.} Per bucket, we compare the mispricing per contract, $\bar\q-\bar\won$
      at the mid, with the cost of trading against it as a taker: the half-spread plus the taker fee
      on the opposite side. The capital term of the band in \cref{sec:framework} is computed at a 5\%
      annual cost of capital over the lead.
\item \emph{Tick floor.} Prices move on a 1\cents{} grid; some markets use 0.1\cents{} steps. On the
      1\cents{} grid, a contract whose fair value is below $0.5\cents$ cannot have a lower mid, and a
      taker cannot pay less than 1\cents. In the 1--2\cents{} bucket we split the loss per contract,
      $\bar\q-\bar\won$, into a mechanical part and a remainder. The mechanical part is the loss of the
      quotes at the floor $\q_{\min}$ (a 0.5\cents{} mid with no bid, or a 1\cents{} ask for takers),
      $s\max(\q_{\min}-\bar\won_{f},0)$, where $s$ is their weight share and $\bar\won_{f}$ their win rate.
      The remainder comes from prices above the floor. Markets on the 0.1\cents{} grid serve as a
      comparison.
\end{itemize}

\paragraph{Robustness.}
We repeat the main measures on several variants of the sample:
\begin{itemize}
\item the first and second halves of the sample, split at the median observation time;
\item without quotes whose mid is at or below 2\cents{} or at or above 98\cents{} (the 1--2\cents{} and
      99\cents{} buckets and their mirror images);
\item with quotes at most 24 hours old, and with spreads of at most 5\cents;
\item by the source of the scheduled time;
\item with design weights, the inverse of each event's sampling probability;
\item with standard errors clustered by market, series or month instead of by event.
\end{itemize}
Many cells are reported. We emphasise patterns that hold across groups and in both halves rather
than single cells.
